% !TEX root = ../Preparazione alle olimpiadi.tex

\chapter*{Introduzione}

I Campionati Italiani di Informatica (ex Olimpiadi) sono una competizione rivolta agli studenti delle istituzioni 
scolastiche secondarie di II grado. Giunte ormai alla XXVI edizione, fanno parte del programma di 
valorizzazione delle eccellenze che la Direzione Generale per gli ordinamenti scolastici e la valutazione del 
sistema nazionale di istruzione del Ministero dell’Istruzione promuove e finanzia ogni anno.

La partecipazione è aperta alle classi I-IV di tutte le istituzioni scolastiche secondarie di II grado, 
statali e paritarie, che ritengano di avere studenti con potenziale interesse per l'informatica, 
soprattutto riguardo gli aspetti logici e algoritmici di tale disciplina. La partecipazione è adatta e 
consigliata anche a coloro che ancora non hanno studiato informatica a livello curricolare, per 
consentire agli studenti di interessarsi e avvicinarsi gradualmente alla materia. 

Gli studenti potranno prepararsi alla prima fase tramite le prove degli anni passati disponibili sul sito 
\url{scolastiche.olinfo.it}.

Il regolamento, inclusa la descrizione delle varie fasi è riportato in appendice \ref{appendice:Regolamento2025}.

Si fa notare che durante la prima fase (selezione scolastica), verrà richiesta la soluzione dei problemi utilizzando esclusivamente \textbf{pseudocodice} le cui regole sono riportate in appendice \ref{appendice:pseudocodice}


In questo documento verranno trattati argomenti teorici e pratici al fine di affrontare le olimpiadi di informatica con una degna preparazione.